% GENERATED FILE. Edit canonical lessons, then run npm run generate:books.
% canonical-source-hash: fnv1a64:1c363db42f8162f5
% canonical-lessons: UR-C537-mutabadil, UR-C537-yaksan, UR-C537-evaz, UR-C537-kasauti, UR-C537-muntakhab-karna

\chapter{Two Ways to Go}
\label{ch:ur-two-ways-to-go}
\hlchaptermodality{537}

\begin{chapteropening}
\textbf{By the end of this chapter:} \emph{I can talk about alternatives and what to choose instead.}

Complete the last lesson of chapter 537: I can talk about alternatives and what to choose instead.
\end{chapteropening}

\section[\texorpdfstring{\ur{متبادل} (mutabādil)}{متبادل (mutabādil)}]{\ur{متبادل} (mutabādil) — an alternative}

\label{lesson:UR-C537-mutabadil}



\begin{quote}
Before the new one: say the Urdu for cover (to shelter behind), then the Urdu for an emergency.
\end{quote}

\subsection*{You'll want to know: \ur{متبادل}}
\textbf{\ur{متبادل}} — \emph{mutabādil} — \textquotedblleft{}an alternative\textquotedblright{}.

\begin{quote}
\ur{اگر} \ur{بس} \ur{نہیں} \ur{آئی} \ur{تو} \ur{رکشہ} \ur{ایک} \ur{متبادل} \ur{ہے۔}
\end{quote}

\begin{quote}
\emph{agar bas nahīṅ āī to rikshā ek mutabādil hai.}
\end{quote}

\begin{quote}
If the bus doesn't come, a rickshaw is another option.
\end{quote}

Masculine.

\begin{cousinweb}
\textbf{\ur{متبادل}} (\emph{mutabādil}) is usually traced, through Persian, to Arabic \emph{mutabādil}, \textquotedblleft{}exchanged, interchangeable\textquotedblright{}.
\end{cousinweb}

\begin{grammarlens}[title={What you've built}]
One more word for saying what happens if.
\end{grammarlens}

\subsection*{Your turn}
\begin{itemize}
\raggedright
  \item \emph{Say it:} \emph{mutabādil}
  \item \emph{Say it:} \emph{mutabādil}, once more
  \item \emph{Say it:} \emph{agar bas nahīṅ āī to rikshā ek mutabādil hai.} — If the bus doesn't come, a rickshaw is another option.
  \item \emph{Recall:} say the Urdu for cover (to shelter behind), then the Urdu for an emergency, then say \emph{mutabādil} again
\end{itemize}

\begin{tcolorbox}[breakable,colback=teal!4,colframe=teal!35!black,title={Before you move on}]
What does \ur{متبادل} mean? (\textquotedblleft{}An alternative\textquotedblright{}.) Say it once more.
\end{tcolorbox}
\section[\texorpdfstring{\ur{یکساں} (yaksāṅ)}{یکساں (yaksāṅ)}]{\ur{یکساں} (yaksāṅ) — alike; the same}

\label{lesson:UR-C537-yaksan}



\begin{quote}
Before the new one: say the Urdu for an emergency, then the Urdu for an alternative.
\end{quote}

\subsection*{You'll want to know: \ur{یکساں}}
\textbf{\ur{یکساں}} — \emph{yaksāṅ} — \textquotedblleft{}alike; the same\textquotedblright{}.

\begin{quote}
\ur{اگر} \ur{دونوں} \ur{راستے} \ur{یکساں} \ur{ہیں} \ur{تو} \ur{کوئی} \ur{بھی} \ur{لیجیے۔}
\end{quote}

\begin{quote}
\emph{agar donoṅ rāste yaksāṅ haiṅ to koī bhī lījiye.}
\end{quote}

\begin{quote}
If both routes are the same, take either one.
\end{quote}

\begin{cousinweb}
\textbf{\ur{یکساں}} (\emph{yaksāṅ}) is usually traced to Persian \emph{yaksān}, \textquotedblleft{}alike\textquotedblright{}, built from \emph{yak}, \textquotedblleft{}one\textquotedblright{}, and \emph{sān}, \textquotedblleft{}manner\textquotedblright{}.
\end{cousinweb}

\begin{grammarlens}[title={What you've built}]
One more word for saying what happens if.
\end{grammarlens}

\subsection*{Your turn}
\begin{itemize}
\raggedright
  \item \emph{Say it:} \emph{yaksāṅ}
  \item \emph{Say it:} \emph{yaksāṅ}, once more
  \item \emph{Say it:} \emph{agar donoṅ rāste yaksāṅ haiṅ to koī bhī lījiye.} — If both routes are the same, take either one.
  \item \emph{Recall:} say the Urdu for an emergency, then the Urdu for an alternative, then say \emph{yaksāṅ} again
\end{itemize}

\begin{tcolorbox}[breakable,colback=teal!4,colframe=teal!35!black,title={Before you move on}]
What does \ur{یکساں} mean? (\textquotedblleft{}Alike; the same\textquotedblright{}.) Say it once more.
\end{tcolorbox}
\section[\texorpdfstring{\ur{عوض} (evaz)}{عوض (evaz)}]{\ur{عوض} (evaz) — something in exchange (for)}

\label{lesson:UR-C537-evaz}



\begin{quote}
Before the new one: say the Urdu for an alternative, then the Urdu for alike.
\end{quote}

\subsection*{You'll want to know: \ur{عوض}}
\textbf{\ur{عوض}} — \emph{evaz} — \textquotedblleft{}something in exchange (for)\textquotedblright{}.

\begin{quote}
\ur{اگر} \ur{قلم} \ur{چاہیے} \ur{تو} \ur{اس} \ur{کے} \ur{عوض} \ur{اپنی} \ur{کتاب} \ur{دیجیے۔}
\end{quote}

\begin{quote}
\emph{agar qalam chāhiye to us ke evaz apnī kitāb dījiye.}
\end{quote}

\begin{quote}
If you want the pen, give me your book in exchange.
\end{quote}

Masculine.

\begin{cousinweb}
\textbf{\ur{عوض}} (\emph{evaz}) is usually traced, through Persian, to Arabic \emph{ʿiwaḍ}, \textquotedblleft{}a substitute, an exchange\textquotedblright{}.
\end{cousinweb}

\begin{grammarlens}[title={What you've built}]
One more word for saying what happens if.
\end{grammarlens}

\subsection*{Your turn}
\begin{itemize}
\raggedright
  \item \emph{Say it:} \emph{evaz}
  \item \emph{Say it:} \emph{evaz}, once more
  \item \emph{Say it:} \emph{agar qalam chāhiye to us ke evaz apnī kitāb dījiye.} — If you want the pen, give me your book in exchange.
  \item \emph{Recall:} say the Urdu for an alternative, then the Urdu for alike, then say \emph{evaz} again
\end{itemize}

\begin{tcolorbox}[breakable,colback=teal!4,colframe=teal!35!black,title={Before you move on}]
What does \ur{عوض} mean? (\textquotedblleft{}Something in exchange (for)\textquotedblright{}.) Say it once more.
\end{tcolorbox}
\section[\texorpdfstring{\ur{کسوٹی} (kasauṭī)}{کسوٹی (kasauṭī)}]{\ur{کسوٹی} (kasauṭī) — a touchstone; a real test}

\label{lesson:UR-C537-kasauti}



\begin{quote}
Before the new one: say the Urdu for alike, then the Urdu for something in exchange (for).
\end{quote}

\subsection*{You'll want to know: \ur{کسوٹی}}
\textbf{\ur{کسوٹی}} — \emph{kasauṭī} — \textquotedblleft{}a touchstone; a real test\textquotedblright{}.

\begin{quote}
\ur{اگر} \ur{دوست} \ur{مشکل} \ur{میں} \ur{ساتھ} \ur{رہا} \ur{تو} \ur{یہی} \ur{کسوٹی} \ur{ہے۔}
\end{quote}

\begin{quote}
\emph{agar dost mushkil meṅ sāth rahā to yihī kasauṭī hai.}
\end{quote}

\begin{quote}
If a friend stands by you in trouble, that is the real test.
\end{quote}

Feminine.

\begin{cousinweb}
\textbf{\ur{کسوٹی}} (\emph{kasauṭī}) is inherited from Sanskrit \emph{kaṣapaṭṭikā}, \textquotedblleft{}a touchstone\textquotedblright{}, the stone gold is rubbed on to test it (CDIAL 2973).
\end{cousinweb}

\begin{grammarlens}[title={What you've built}]
One more word for saying what happens if.
\end{grammarlens}

\subsection*{Your turn}
\begin{itemize}
\raggedright
  \item \emph{Say it:} \emph{kasauṭī}
  \item \emph{Say it:} \emph{kasauṭī}, once more
  \item \emph{Say it:} \emph{agar dost mushkil meṅ sāth rahā to yihī kasauṭī hai.} — If a friend stands by you in trouble, that is the real test.
  \item \emph{Recall:} say the Urdu for alike, then the Urdu for something in exchange (for), then say \emph{kasauṭī} again
\end{itemize}

\begin{tcolorbox}[breakable,colback=teal!4,colframe=teal!35!black,title={Before you move on}]
What does \ur{کسوٹی} mean? (\textquotedblleft{}A touchstone; a real test\textquotedblright{}.) Say it once more.
\end{tcolorbox}
\section[\texorpdfstring{\ur{منتخب} \ur{کرنا} (muntakhab karnā)}{منتخب کرنا (muntakhab karnā)}]{\ur{منتخب} \ur{کرنا} (muntakhab karnā) — to select (formally)}

\label{lesson:UR-C537-muntakhab-karna}



\begin{quote}
Before the new one: say the Urdu for something in exchange (for), then the Urdu for a touchstone.
\end{quote}

\subsection*{You'll want to know: \ur{منتخب} \ur{کرنا}}
\textbf{\ur{منتخب} \ur{کرنا}} — \emph{muntakhab karnā} — \textquotedblleft{}to select (formally)\textquotedblright{}.

\begin{quote}
\ur{اگر} \ur{میں} \ur{نے} \ur{یہ} \ur{کتاب} \ur{منتخب} \ur{کی} \ur{تو} \ur{پڑھوں} \ur{گی۔}
\end{quote}

\begin{quote}
\emph{agar maiṅ ne yih kitāb muntakhab kī to paṛhūṅ gī.}
\end{quote}

\begin{quote}
If I pick this book, I will read it.
\end{quote}

Past, with \ur{نے} (he or she …): \ur{اس} \ur{نے} \ur{ایک} \ur{کتاب} \ur{منتخب} \ur{کی}, \emph{us ne ek kitāb muntakhab kī}.

\begin{cousinweb}
\textbf{\ur{منتخب} \ur{کرنا}} (\emph{muntakhab karnā}) is built from \emph{muntakhab}, \textquotedblleft{}chosen\textquotedblright{}, usually traced, through Persian, to Arabic \emph{muntakhab}, and \emph{karnā}, \textquotedblleft{}to do\textquotedblright{}.
\end{cousinweb}

\begin{grammarlens}[title={What you've built}]
One more word for saying what happens if.
\end{grammarlens}

\subsection*{Your turn}
\begin{itemize}
\raggedright
  \item \emph{Say it:} \emph{muntakhab karnā}
  \item \emph{Say it:} \emph{muntakhab karnā}, once more
  \item \emph{Say it:} \emph{agar maiṅ ne yih kitāb muntakhab kī to paṛhūṅ gī.} — If I pick this book, I will read it.
  \item \emph{Recall:} say the Urdu for something in exchange (for), then the Urdu for a touchstone, then say \emph{muntakhab karnā} again
\end{itemize}

\begin{tcolorbox}[breakable,colback=teal!4,colframe=teal!35!black,title={Before you move on}]
What does \ur{منتخب} \ur{کرنا} mean? (\textquotedblleft{}To select (formally)\textquotedblright{}.) Say it once more.
\end{tcolorbox}
